%% ============================================================
%%  Chapter 1 — Introduction
%%  01_introduction.tex
%% ============================================================
\chapter{Introduction}

\chapterintro{%
  This chapter establishes the full business and technical context of the
  Clinic Management API. It covers the executive summary, the technology
  stack, the RESTful design philosophy, and an exhaustive description of
  every entity within the core business domain. Readers who are already
  familiar with the project may skip to Chapter~\ref{chap:architecture}
  for the system architecture breakdown.
}

%% ─────────────────────────────────────────────────────────────────────────
\section{Executive Summary}
%% ─────────────────────────────────────────────────────────────────────────

The Clinic Management API is an enterprise-grade, RESTful backend system
built on the Laravel framework. It is designed to serve as the single
authoritative data layer for one or more medical clinic operations. The
system digitises and enforces the complete patient journey: from first
registration, through appointment scheduling and confirmation, to the
issuance of a formal doctor's prescription.

At its core, the API solves a fundamental coordination problem present in
any multi-stakeholder clinical environment. Without a unified system,
scheduling is handled via telephone or paper registers, patient records
are stored in disparate file systems, and prescriptions are written by
hand with no linkage to the underlying appointment. These fragmented
workflows introduce avoidable medical errors, administrative overhead,
and a complete absence of an auditable data trail.

The Clinic Management API replaces all of these manual processes with a
structured, role-aware, and cryptographically authenticated JSON API.
Every action --- listing available appointment slots, booking a
consultation, updating a patient's record, or issuing a multi-item
prescription --- is mediated by explicit business rules enforced at the
service layer, not merely at the database constraint level. This
distinction is deliberate and architecturally significant: business rules
that live in the database are invisible to the application; business rules
that live in the service layer are testable, versionable, and auditable.

The system supports three discrete user roles: \textbf{Super Administrator},
\textbf{Doctor}, and \textbf{Assistant}. Each role carries a precisely
defined permission surface. A Doctor may confirm or complete appointments
that belong to them but cannot create new clinic records or modify another
doctor's prescriptions. An Assistant may register new patients and book
appointments on a patient's behalf but cannot alter appointment status
beyond cancellation or issue prescriptions. A Super Administrator operates
globally across all clinics and is the only actor authorised to manage
the user roster itself.

The API is designed to be stateless, horizontally scalable, and fully
decoupled from any frontend technology. A mobile application, a web
dashboard, and a third-party electronic health record system can all
consume the same endpoints simultaneously, each receiving the same
consistently structured JSON envelope.

%% ─────────────────────────────────────────────────────────────────────────
\section{Purpose and Scope of This Document}
%% ─────────────────────────────────────────────────────────────────────────

This technical documentation manual serves four distinct audiences, each
with different information requirements.

\begin{enumerate}
  \item \textbf{Backend Engineers} maintaining or extending the API need a
    precise reference for the repository-service-controller layering, the
    custom exception hierarchy, and the business rules encoded in each
    service class.

  \item \textbf{QA Engineers} writing or reviewing the PHPUnit feature test
    suite need a complete mapping from business requirement to test case,
    including the exact HTTP method, route, request payload, expected
    response envelope, and HTTP status code for every scenario.

  \item \textbf{Frontend and Mobile Developers} integrating against the API
    need the endpoint catalogue in Chapter~\ref{chap:endpoints}, the
    authentication flow in Chapter~\ref{chap:auth}, and the canonical JSON
    response envelope format described in Section~\ref{sec:response-envelope}.

  \item \textbf{Technical Architects} evaluating the system need the
    architectural diagrams, the database entity-relationship analysis, and
    the exception handling strategy to assess compliance with their
    organisation's engineering standards.
\end{enumerate}

The scope of this document is limited to the backend API layer. It does
not cover any specific frontend or mobile client implementation, CI/CD
pipeline configuration, server provisioning, or HTTPS termination
strategy, all of which are treated as deployment concerns external to the
application itself.

The document is structured to be read sequentially on first pass and
used as a reference thereafter. Each chapter is self-contained, with
cross-references provided wherever concepts from another chapter are
relevant.

%% ─────────────────────────────────────────────────────────────────────────
\section{Document Structure Overview}
%% ─────────────────────────────────────────────────────────────────────────

\begin{table}[h!]
  \centering
  \caption{Chapter map and primary audience for each section.}
  \label{tab:chapter-map}
  \renewcommand{\arraystretch}{1.4}
  \begin{tabularx}{\linewidth}{@{}C{0.5cm} L{4.5cm} X L{3.2cm}@{}}
    \toprule
    \rowcolor{PrimaryBlue}
    \tblhdr{\#} & \tblhdr{Chapter} & \tblhdr{Content Summary} & \tblhdr{Primary Audience} \\
    \midrule
    1 & Introduction          & Business domain, tech stack, conventions          & All \\
    \rowalt
    2 & System Architecture   & Layered design, middleware pipeline, directory
                                structure                                          & Engineers, Architects \\
    3 & Database \& Models    & ERD, migrations, Eloquent models, relationships,
                                UUID strategy                                     & Engineers \\
    \rowalt
    4 & Authentication \& Roles & Passport OAuth2 flow, RBAC, middleware chain,
                                  clinic scope enforcement                        & Engineers, Frontend \\
    5 & API Endpoints         & Full endpoint catalogue with request/response
                                examples for every route                          & All \\
    \rowalt
    6 & Exception Handling    & Custom exception hierarchy, HTTP status mapping,
                                bootstrap registration                            & Engineers, QA \\
    7 & Testing Strategy      & PHPUnit configuration, factory design, base test
                                class, seeding strategy                           & QA, Engineers \\
    \rowalt
    8 & Test Cases \& Validation & Complete test case matrix: happy path (Phase~1)
                                   and alternative scenarios (Phase~2)            & QA \\
    \bottomrule
  \end{tabularx}
\end{table}

%% ─────────────────────────────────────────────────────────────────────────
\section{Technology Stack}
\label{sec:tech-stack}
%% ─────────────────────────────────────────────────────────────────────────

The Clinic Management API is built on a tightly integrated, open-source
technology stack. Every component was chosen for its maturity, its
community support, its security track record, and its first-class
integration with the Laravel ecosystem.

%% ── §4.1 ─────────────────────────────────────────────────────────────────
\subsection{Runtime Environment: PHP 8.4}

PHP 8.4 is the minimum required runtime version for this application. The
choice of PHP 8.4 over earlier minor versions is not arbitrary; it
unlocks several language features that are used pervasively throughout
the codebase.

\textbf{Named arguments} allow service-class constructors and factory
state methods to be called with keyword syntax, improving readability in
deeply nested calls without requiring intermediate variables.

\textbf{Union types} and \textbf{intersection types} allow method
signatures to be precise about which concrete types they accept and
return, eliminating an entire category of runtime type errors that would
otherwise only be discovered in production.

\textbf{Readonly properties} (introduced in PHP 8.1 and refined in 8.2
and beyond) are used in value objects and data-transfer objects to
guarantee that state cannot be mutated after construction, which enforces
immutability at the language level rather than by convention.

\textbf{Fibers} (PHP 8.1) provide the coroutine primitive that underpins
Laravel's asynchronous queue and job infrastructure, which is leveraged by
the appointment confirmation job system described in
Section~\ref{sec:jobs}.

\textbf{Enums} (PHP 8.1) are the mechanism by which the application
models \texttt{UserRole} and \texttt{AppointmentStatus} as first-class
types rather than magic strings. This prevents an entire class of
runtime errors caused by mis-typed status values and integrates cleanly
with Eloquent's attribute casting system.

\begin{table}[h!]
  \centering
  \caption{PHP features used in the Clinic Management API.}
  \label{tab:php-features}
  \renewcommand{\arraystretch}{1.4}
  \begin{tabularx}{\linewidth}{@{}L{3.8cm} C{2.2cm} X@{}}
    \toprule
    \rowcolor{PrimaryBlue}
    \tblhdr{Feature} & \tblhdr{Introduced} & \tblhdr{Usage in This Project} \\
    \midrule
    Enums                    & PHP 8.1 & \texttt{UserRole}, \texttt{AppointmentStatus},
                                         \texttt{Gender} \\
    \rowalt
    Readonly Properties      & PHP 8.1 & DTO and value objects \\
    Named Arguments          & PHP 8.0 & Factory state methods, service constructors \\
    \rowalt
    Union \& Nullable Types  & PHP 8.0 & Repository method signatures \\
    Match Expressions        & PHP 8.0 & Status transition guards \\
    \rowalt
    Constructor Promotion    & PHP 8.0 & All service and repository classes \\
    Fibers                   & PHP 8.1 & Queue job infrastructure \\
    \rowalt
    First-class Callables    & PHP 8.1 & Collection \texttt{map()} and
                                         \texttt{filter()} chains \\
    \bottomrule
  \end{tabularx}
\end{table}

%% ── §4.2 ─────────────────────────────────────────────────────────────────
\subsection{Web Framework: Laravel 12}

Laravel 12 is the application framework. Laravel provides the HTTP
routing layer, the dependency injection container, the Eloquent ORM, the
queue system, the cache abstraction, and the testing harness. The version
12 release series is built on top of the Symfony HTTP Foundation component,
ensuring standards compliance with PSR-7 request/response semantics while
adding the rich developer experience that the Laravel ecosystem is known for.

The key Laravel subsystems in active use are as follows.

\textbf{Routing.} All API routes are defined in
\texttt{routes/api.php} and are prefixed with \texttt{/api/\apiversion{}/}
by the framework's route service provider. Route model binding is used
throughout the application to automatically resolve Eloquent model
instances from UUID path segments, eliminating repetitive
\texttt{findOrFail} boilerplate from controller methods.

\textbf{Middleware Pipeline.} Every request passes through a layered
middleware stack before reaching its controller. The \texttt{auth:api}
middleware (provided by Laravel Passport) authenticates the bearer token.
The custom \texttt{RoleMiddleware} enforces role-based access control.
The custom \texttt{ClinicScopeMiddleware} injects the authenticated
user's \texttt{clinic\_id} into the request data bag, preventing doctors
and assistants from supplying a fabricated clinic identifier in their
JSON payload.

\textbf{Form Request Validation.} Dedicated \texttt{FormRequest} classes
are used for every mutating endpoint. This separates validation logic
from controller logic and provides automatic 422 responses with
structured field-level error messages when validation fails, without a
single explicit \texttt{if} statement in any controller.

\textbf{Eloquent ORM.} Laravel's active-record ORM is used for all
database interactions. Models declare their relationships (one-to-many,
belongs-to) as methods, enabling eager loading of related data to
avoid N+1 query problems. The \texttt{HasUuids} trait is applied to
every model to ensure that primary keys are generated as UUIDs rather
than auto-incrementing integers.

\textbf{Service Container.} All services, repositories, and third-party
clients are resolved through Laravel's dependency injection container.
Constructor injection is used universally, which enables clean unit
testing through mock substitution and ensures that no service class
instantiates its own dependencies.

\textbf{Cache Abstraction.} The appointment slot availability system
uses Laravel's cache facade with a Redis or file driver (configurable via
\texttt{CACHE\_DRIVER}) to memoise the result of expensive slot computation
for a five-minute window. A distributed cache lock (\texttt{Cache::lock})
is used during the booking write path to prevent race conditions that
could result in two patients booking the same time slot simultaneously.

\textbf{Queue System.} The \texttt{SendAppointmentConfirmation} job is
dispatched asynchronously after every successful booking. The queue
driver is configurable and defaults to the synchronous driver in the
testing environment to prevent test isolation issues.

%% ── §4.3 ─────────────────────────────────────────────────────────────────
\subsection{Database: MySQL 8.0}

MySQL 8.0 is the sole supported database engine for this application.
This is not a portability oversight; it is a deliberate design decision
that enables the use of MySQL-specific features that provide measurable
value in a production clinical environment.

\textbf{FULLTEXT Indexes.} The \texttt{patients} table has a FULLTEXT
index defined across the \texttt{name} and \texttt{phone} columns. The
patient search endpoint uses \texttt{whereFullText} to perform
natural-language full-text search, which is significantly more
performant than \texttt{LIKE '\%...\%'} queries on large patient rosters
and is a MySQL-specific capability not available in SQLite.

\textbf{UUID Primary Keys.} All primary keys are stored as
\texttt{CHAR(36)} columns containing RFC~4122 UUID v4 values generated
by PHP's \texttt{Str::uuid()} method via the Eloquent
\texttt{HasUuids} trait. This choice has several consequences. UUIDs are
safe to expose in API responses because they reveal nothing about the
total number of records, the insertion order, or the internal
structure of the database. They also enable records to be created in a
distributed environment (multiple application servers, offline-first
mobile clients) without a centralised auto-increment authority.

\textbf{TIME Column Type.} Doctor schedule start and end times and
appointment slot times are stored as MySQL's native \texttt{TIME} type.
MySQL returns TIME values over the wire as strings formatted as
\texttt{HH:MM:SS}. The application normalises these to \texttt{HH:MM}
format in the repository layer before business logic comparisons, a
subtlety that is documented in detail in
Section~\ref{sec:time-normalisation}.

\textbf{Unique Composite Indexes.} The \texttt{appointments} table
carries a unique index on the composite of
\texttt{(doctor\_id, appointment\_date, start\_time)}, enforced at the
database level as a last line of defence against double-booking, even if
the application-layer cache-lock strategy is somehow bypassed.

\textbf{Referential Integrity.} All foreign key relationships are
enforced with \texttt{FOREIGN KEY} constraints and
\texttt{ON DELETE CASCADE} or \texttt{ON DELETE RESTRICT} semantics as
appropriate to prevent orphaned records.

%% ── §4.4 ─────────────────────────────────────────────────────────────────
\subsection{Authentication: Laravel Passport (OAuth2)}

Authentication is handled exclusively by Laravel Passport, which
implements the full OAuth2 specification on top of the application's
user model. The API uses the \textbf{Personal Access Token} (PAT) grant
type, which is the appropriate OAuth2 grant for a first-party API where
the client and the server are the same organisation.

When a user authenticates via the \texttt{POST /api/\apiversion{}/auth/login}
endpoint, the \texttt{AuthService} verifies the provided credentials
against the hashed password in the database, checks the
\texttt{is\_active} flag, and then calls \texttt{createToken()} on the
Eloquent user model. Passport generates a cryptographically signed
access token and persists a token record to its \texttt{oauth\_access\_tokens}
table. The raw token string is returned to the client only once; it is
never stored in plaintext on the server.

On every subsequent request, the client presents the token in the HTTP
\texttt{Authorization} header as a Bearer token. The \texttt{auth:api}
Passport middleware validates the token's signature, checks whether the
token has been revoked, and populates \texttt{request()->user()} with
the authenticated Eloquent model instance. All of this occurs before the
request enters the application's own middleware pipeline.

Token revocation (logout) is achieved by calling \texttt{revoke()} on
the token instance, which sets the \texttt{revoked} flag in the
\texttt{oauth\_access\_tokens} table. Subsequent requests bearing the
same token will be rejected with a \stunauth{} response.

%% ── §4.5 ─────────────────────────────────────────────────────────────────
\subsection{Testing Framework: PHPUnit 11}

The application ships with a comprehensive PHPUnit 11 feature test suite
comprising \testcount{} passing tests and 254 assertions. All tests are
\textbf{feature tests} that exercise the full HTTP stack from the HTTP
request layer down to the MySQL database, ensuring that the integration
between the routing layer, middleware, controllers, services, repositories,
and Eloquent models is continuously verified.

The test suite uses the \texttt{RefreshDatabase} trait, which wraps
each test in a database transaction that is rolled back after the test
completes, leaving the database in a clean state for the next test.
Because the application is configured to use a real MySQL database
(not an in-memory SQLite substitute), the tests validate MySQL-specific
behaviour such as FULLTEXT search and TIME column formatting.

Test data is generated exclusively through Laravel model factories, with
dedicated factories for every model in the system. Factory state methods
(\texttt{superAdmin()}, \texttt{doctor()}, \texttt{assistant()},
\texttt{forClinic()}, \texttt{inactive()}) provide a clean, composable
DSL for expressing complex test data scenarios.

\begin{table}[h!]
  \centering
  \caption{Full technology stack summary.}
  \label{tab:tech-stack}
  \renewcommand{\arraystretch}{1.5}
  \begin{tabularx}{\linewidth}{@{}L{3.8cm} L{3.2cm} X@{}}
    \toprule
    \rowcolor{PrimaryBlue}
    \tblhdr{Component} & \tblhdr{Technology} & \tblhdr{Version / Notes} \\
    \midrule
    Runtime            & PHP          & 8.4+ required \\
    \rowalt
    Framework          & Laravel      & 12.x \\
    Database           & MySQL        & 8.0+ (FULLTEXT, UUID, TIME) \\
    \rowalt
    Authentication     & Laravel Passport & OAuth2 Personal Access Tokens \\
    Caching            & Laravel Cache & Redis recommended; file driver for dev \\
    \rowalt
    Queue System       & Laravel Queue & sync driver for tests; Redis for prod \\
    Testing            & PHPUnit      & 11.x --- \testcount{} tests, 254 assertions \\
    \rowalt
    API Protocol       & REST over HTTP/1.1 & JSON request and response bodies \\
    Primary Keys       & UUID v4      & \texttt{HasUuids} on every Eloquent model \\
    \rowalt
    Deployment Target  & Linux / Docker & PHP-FPM + Nginx recommended \\
    \bottomrule
  \end{tabularx}
\end{table}

%% ─────────────────────────────────────────────────────────────────────────
\section{RESTful API Design Philosophy}
\label{sec:rest-philosophy}
%% ─────────────────────────────────────────────────────────────────────────

The API adheres to the REST architectural style as defined by Roy
Fielding in his 2000 doctoral dissertation. The six Fielding constraints
--- client-server separation, statelessness, cacheability, uniform
interface, layered system, and code-on-demand (not used) --- are
observed throughout the design.

\subsection{Statelessness}

Each HTTP request contains all the information necessary for the server
to fulfil it. The server retains no session state between requests. The
authentication token, the user's identity, and the user's clinic
affiliation are all derived entirely from the bearer token on each
request. This makes the application trivially horizontally scalable:
any application server instance can handle any request because no
session affinity is required.

\subsection{Uniform Interface and Resource Naming}

Resources are named as lowercase plural nouns in the URL path. Actions
are expressed through the HTTP method rather than through the URL. The
base path for all API resources is \route{/api/\apiversion{}/}, where
\texttt{\apiversion{}} is the version string that allows the API to
evolve without breaking existing clients.

\begin{table}[h!]
  \centering
  \caption{RESTful URL and method conventions used in this API.}
  \label{tab:rest-conventions}
  \renewcommand{\arraystretch}{1.5}
  \begin{tabularx}{\linewidth}{@{}L{2cm} L{6cm} X@{}}
    \toprule
    \rowcolor{PrimaryBlue}
    \tblhdr{Method} & \tblhdr{URL Pattern} & \tblhdr{Semantic} \\
    \midrule
    \mget    & \route{/resources}          & List a paginated collection \\
    \rowalt
    \mpost   & \route{/resources}          & Create a new resource \\
    \mget    & \route{/resources/\{id\}}   & Retrieve a single resource by ID \\
    \rowalt
    \mput    & \route{/resources/\{id\}}   & Replace / update a resource \\
    \mpatch  & \route{/resources/\{id\}/status} & Partial update of a sub-field \\
    \rowalt
    \mdelete & \route{/resources/\{id\}}   & Delete or cancel a resource \\
    \bottomrule
  \end{tabularx}
\end{table}

\subsection{JSON Response Envelope}
\label{sec:response-envelope}

Every response from the API --- success or failure --- is wrapped in a
consistent JSON envelope. This envelope provides a predictable structure
that client code can rely on regardless of the endpoint being called.

\begin{lstlisting}[style=json, caption={Standard success response envelope.}]
{
  "success": true,
  "message": "Human-readable description of the outcome",
  "data": { }
}
\end{lstlisting}

\begin{lstlisting}[style=json, caption={Standard error response envelope.}]
{
  "success": false,
  "message": "Human-readable description of the error"
}
\end{lstlisting}

\begin{lstlisting}[style=json, caption={Validation failure envelope (HTTP 422).}]
{
  "success": false,
  "errors": {
    "field_name": ["The field_name field is required."]
  }
}
\end{lstlisting}

The \texttt{success} boolean allows client code to branch on a single
field check rather than inspecting the HTTP status code. The
\texttt{message} string is always human-readable and suitable for display
in a user interface notification. The \texttt{data} field is present
only on successful responses and contains either a single resource
object or a flat array of resource objects.

\begin{notebox}[Pagination]
  List endpoints return a flat array in the \texttt{data} field. The API
  uses Laravel's \texttt{LengthAwarePaginator} internally, but the
  \texttt{AnonymousResourceCollection::toArray()} call in the
  \texttt{ApiResponse} trait resolves the paginated result into a flat
  array before returning. There is no \texttt{meta} or \texttt{links}
  key at the top level of the envelope. Pagination parameters
  (\texttt{page}, \texttt{per\_page}) are accepted as query string
  parameters.
\end{notebox}

\subsection{Versioning Strategy}

The API is versioned at the URL path level. All routes are nested under
the prefix \route{/api/v1/}. This prefix is registered in
\texttt{routes/api.php} and enforced by the routing configuration in
\texttt{bootstrap/app.php}. When breaking changes must be introduced in
the future, a new route group under \route{/api/v2/} can be registered
without disturbing existing \texttt{v1} clients.

\subsection{HTTP Status Code Discipline}

The API uses HTTP status codes as primary signalling mechanisms, not as
secondary annotations on a separate success field. The following status
codes are used with strict semantic meaning.

\begin{table}[h!]
  \centering
  \caption{HTTP status codes and their semantic meaning in this API.}
  \label{tab:status-codes}
  \renewcommand{\arraystretch}{1.5}
  \begin{tabularx}{\linewidth}{@{}C{2.5cm} X@{}}
    \toprule
    \rowcolor{PrimaryBlue}
    \tblhdr{Status Code} & \tblhdr{Meaning and Trigger} \\
    \midrule
    \stok         & Request succeeded; existing resource returned in \texttt{data}. \\
    \rowalt
    \stcreated    & Request succeeded; new resource was created and returned in
                    \texttt{data}. \\
    \stunauth     & No valid bearer token was presented; the
                    \texttt{auth:api} middleware rejected the request. \\
    \rowalt
    \stforbid     & The authenticated user does not have the required role
                    (\texttt{RoleMiddleware}), or is attempting to access a
                    resource outside their clinic's scope
                    (\texttt{ClinicScopeViolationException}), or is triggering
                    a self-demotion (\texttt{SelfDemotionException}), or their
                    account is deactivated (\texttt{AccountDeactivatedException}). \\
    \stnotfound   & Route model binding failed to resolve the UUID in the path
                    to an existing Eloquent model instance. \\
    \rowalt
    \stconflict   & A booking or reschedule request targets a time slot that is
                    already occupied (\texttt{AppointmentConflictException}). \\
    \stunproc     & A form request validation failed, or an operation was
                    attempted on a resource in an incompatible state
                    (\texttt{InvalidAppointmentStateException}). \\
    \bottomrule
  \end{tabularx}
\end{table}

%% ─────────────────────────────────────────────────────────────────────────
\section{Core Business Domain}
\label{sec:business-domain}
%% ─────────────────────────────────────────────────────────────────────────

The application models a small-to-medium medical clinic chain. The
domain entities and their relationships are described in full in this
section. A formal entity-relationship diagram and the corresponding
database schema are provided in Chapter~\ref{chap:database}, but the
narrative description here is intended to give engineers an intuitive
understanding of the system's intent before they encounter the
technical schema.

%% ── §6.1 ─────────────────────────────────────────────────────────────────
\subsection{The Clinic Entity}
\label{sec:domain-clinic}

A \textbf{Clinic} is the fundamental organisational unit of the system.
Every other entity --- users, patients, appointments, and prescriptions
--- belongs to exactly one clinic. The clinic boundary is the primary
data isolation boundary in the application: no actor operating within
one clinic can read or write data belonging to a different clinic,
except for the Super Administrator, who operates at the global level
above all clinic boundaries.

A clinic has the following attributes:

\begin{itemize}
  \item \textbf{Name} --- the legal or trading name of the clinic,
    unique across the system to prevent confusion in administrative views.
  \item \textbf{Email} --- a contact email address for the clinic,
    unique across the system. This is used for external communications
    such as appointment confirmation emails.
  \item \textbf{Address} --- the physical street address of the clinic.
  \item \textbf{Phone} --- the primary telephone number.
  \item \textbf{Is Active} --- a boolean flag. A deactivated clinic
    (\texttt{is\_active = false}) is retained in the database for
    historical integrity but its users and patients cannot be accessed
    through normal API workflows. This supports soft-discontinuation of
    a clinic branch without destroying its historical appointment and
    prescription records.
\end{itemize}

Clinics are managed exclusively by the Super Administrator. No Doctor or
Assistant has any create, update, or delete capability over clinic records.

%% ── §6.2 ─────────────────────────────────────────────────────────────────
\subsection{The User and Role System}
\label{sec:domain-users}

A \textbf{User} represents a staff member of the clinic chain. The
system does not model patients as users; patients are a separate,
distinct entity (see Section~\ref{sec:domain-patients}). Every user has
an assigned role, modelled as a PHP 8.1 backed enum:
\texttt{UserRole::SUPER\_ADMIN}, \texttt{UserRole::DOCTOR}, or
\texttt{UserRole::ASSISTANT}.

Every user (except the Super Administrator) is bound to a single clinic
via the \texttt{clinic\_id} foreign key. A user cannot be transferred
from one clinic to another without the Super Administrator explicitly
updating their \texttt{clinic\_id}. This single-clinic constraint is
the foundation of the clinic-scope enforcement described in
Section~\ref{sec:clinic-scope}.

\subsubsection{Super Administrator}

The Super Administrator (\texttt{UserRole::SUPER\_ADMIN}) is the global
administrator of the entire system. There is no architectural restriction
on the number of Super Administrator accounts, but in practice a single
organisation typically maintains one or two.

The Super Administrator's permission surface covers:

\begin{itemize}
  \item Full CRUD over all \textbf{Clinic} records.
  \item Full CRUD over all \textbf{User} records, including the ability
    to assign or change roles and to activate or deactivate accounts.
  \item Read access to the \textbf{Dashboard} aggregate, which returns
    system-wide KPIs: total clinics, total users, total patients, total
    appointments, and a date-range-filterable appointment breakdown.
  \item The Super Administrator is the only actor who can call the
    \texttt{PATCH /api/\apiversion{}/admin/users/\{id\}/role} endpoint
    to change a user's role. A hard-coded guard in the
    \texttt{UserController} prevents a Super Administrator from
    demoting their own account via this endpoint; the system throws a
    \phpclass{SelfDemotionException} if this is attempted.
\end{itemize}

The Super Administrator does \emph{not} have access to the Doctor or
Assistant route groups. The role middleware enforces strict, non-overlapping
role surfaces. A Super Administrator who needs to view clinical data
must be assigned a secondary Doctor or Assistant account, or a future
scope-expansion endpoint must be explicitly added.

\subsubsection{Doctor}

The Doctor (\texttt{UserRole::DOCTOR}) is a licensed medical practitioner
affiliated with exactly one clinic. The Doctor's permission surface
covers all clinically sensitive data within their own clinic.

\begin{itemize}
  \item \textbf{Schedule Management.} A Doctor defines their own weekly
    availability through the \texttt{DoctorSchedule} entity. A schedule
    entry specifies a day of the week (0 = Sunday through 6 = Saturday),
    a start time, an end time, and a slot duration in minutes. The
    appointment slot generator uses this schedule to compute the list
    of bookable time slots for any given date.
  \item \textbf{Appointment Visibility.} A Doctor can list and view all
    appointments where they are the assigned doctor. They can filter
    by date and status. They cannot see appointments assigned to other
    doctors, even within the same clinic.
  \item \textbf{Appointment Status Management.} A Doctor may advance
    the status of their own appointments from \texttt{pending} to
    \texttt{confirmed}, from \texttt{confirmed} to \texttt{completed}.
    Cancellation is handled by the Assistant. There is no direct
    Doctor cancellation endpoint in the current scope.
  \item \textbf{Prescription Authoring.} A Doctor is the only role that
    may create a prescription. Prescriptions are always linked to a
    specific appointment, and the system enforces that only the
    doctor assigned to that appointment may author its prescription.
    Prescriptions support multiple line items (medicines), each with
    a dosage, frequency, duration, and optional notes.
  \item \textbf{Patient Read Access.} A Doctor may list and view patients
    within their clinic. They do not have write access to patient records
    (patient registration is the Assistant's responsibility). Access to
    a patient from another clinic raises a
    \phpclass{ClinicScopeViolationException}.
\end{itemize}

\subsubsection{Assistant}

The Assistant (\texttt{UserRole::ASSISTANT}) is the front-desk or
administrative staff member. The Assistant's role is to manage patient
intake and appointment logistics on behalf of the clinic and its doctors.

\begin{itemize}
  \item \textbf{Patient Management.} The Assistant has full create and
    update access to patient records within their clinic. They can
    register a new patient (providing name, email, phone, date of birth,
    gender, and address), search the patient roster, and update patient
    details. The clinic scope is automatically enforced: the
    \texttt{ClinicScopeMiddleware} injects the assistant's
    \texttt{clinic\_id} into the request, preventing cross-clinic
    patient creation or modification.
  \item \textbf{Appointment Booking.} The Assistant may book an
    appointment on behalf of a patient by supplying the doctor's ID,
    the patient's ID, a date, and a start time. The system validates
    that the chosen slot is available (not already booked and within
    the doctor's defined schedule for that day of the week) and
    computes the end time automatically from the slot duration. A
    distributed cache lock prevents two concurrent bookings from
    claiming the same slot.
  \item \textbf{Appointment Rescheduling.} The Assistant may reschedule
    a pending or confirmed appointment by providing a new date and
    start time. The same slot availability validation applies. A
    completed appointment cannot be rescheduled; the system throws an
    \phpclass{InvalidAppointmentStateException}.
  \item \textbf{Appointment Cancellation.} The Assistant may cancel
    any appointment within their clinic by calling the
    \texttt{DELETE} endpoint. The appointment is not removed from the
    database; its status is updated to \texttt{cancelled} and the
    time slot is released back into the availability pool.
  \item \textbf{Available Slot Enquiry.} The Assistant may query the
    list of available time slots for a specific doctor on a specific
    date, which drives the appointment booking UI without exposing raw
    schedule data.
\end{itemize}

%% ── §6.3 ─────────────────────────────────────────────────────────────────
\subsection{The Patient Entity}
\label{sec:domain-patients}

A \textbf{Patient} is a person who receives care at a clinic. Patients
are not system users and do not authenticate with the API. They exist
purely as data objects that are created and maintained by Assistant-role
users.

\begin{itemize}
  \item \textbf{Name} --- Full legal name. Indexed with FULLTEXT for
    search capability.
  \item \textbf{Phone} --- Primary contact number. Indexed with FULLTEXT
    alongside name to support search by phone number fragment. Must be
    unique within the clinic to prevent duplicate patient registration.
  \item \textbf{Email} --- Optional email address for communication.
  \item \textbf{Date of Birth} --- Used in clinical contexts for age
    verification and dosage calculation.
  \item \textbf{Gender} --- Modelled as the PHP enum
    \texttt{Gender::MALE} or \texttt{Gender::FEMALE}, cast via
    Eloquent's enum casting. Returned in API responses as the lowercase
    string value (\texttt{"male"} or \texttt{"female"}).
  \item \textbf{Address} --- Optional residential address.
  \item \textbf{Clinic ID} --- The foreign key binding the patient to
    their clinic. Set automatically by the \texttt{ClinicScopeMiddleware}
    and cannot be overridden by the client.
\end{itemize}

\begin{warningbox}[Cross-Clinic Patient Access]
  It is a deliberate security constraint that a patient record belonging
  to \emph{Clinic A} cannot be read, updated, or referenced in an
  appointment by any user of \emph{Clinic B}. Attempting to do so will
  result in a \stforbid{} response via the
  \phpclass{ClinicScopeViolationException}.
\end{warningbox}

%% ── §6.4 ─────────────────────────────────────────────────────────────────
\subsection{Doctor Schedules and Time Slot Management}
\label{sec:domain-schedules}

A \textbf{DoctorSchedule} record defines the recurring weekly
availability of a specific doctor. Each record encodes the following
information:

\begin{itemize}
  \item \textbf{Doctor ID} --- the UUID of the doctor this schedule
    belongs to.
  \item \textbf{Clinic ID} --- the clinic the schedule is associated
    with, used for scope validation.
  \item \textbf{Day of Week} --- an integer from 0 (Sunday) to 6
    (Saturday) following the PHP \texttt{Carbon::dayOfWeek} convention.
  \item \textbf{Start Time} and \textbf{End Time} --- the beginning and
    end of the doctor's working block for the given day.
  \item \textbf{Slot Duration} --- the length of each appointment slot,
    in minutes. A slot duration of 60 minutes means the doctor takes
    one patient per hour.
\end{itemize}

A doctor may have multiple schedule entries (one per working day), but
at most one schedule per day per doctor. The unique constraint on
\texttt{(doctor\_id, day\_of\_week)} in the database enforces this.

The \textbf{slot generation algorithm} in \texttt{AppointmentService}
operates as follows: given a doctor's ID and a target date, the service
looks up the schedule for the corresponding day of the week. It then
iterates from the schedule's start time to its end time in increments
of \texttt{slot\_duration} minutes, generating a list of candidate
start times. From this list, it subtracts any start times that already
have a non-cancelled appointment record. The resulting array is the set
of available slots returned to the client.

\begin{tipbox}[Time Format Normalisation]
  MySQL stores TIME columns and returns them over the wire as
  \texttt{HH:MM:SS} (e.g.\ \texttt{08:00:00}). The slot generation
  algorithm formats candidate slots as \texttt{H:i} (e.g.\
  \texttt{08:00}). Without normalisation, the \texttt{in\_array}
  check would always fail, making every slot appear available even
  when booked. The \texttt{AppointmentRepository::getBookedTimeSlots()}
  method normalises MySQL's output using a \texttt{substr(\$t, 0, 5)}
  call before returning the array, ensuring format consistency.
\end{tipbox}

%% ── §6.5 ─────────────────────────────────────────────────────────────────
\subsection{The Appointment Lifecycle}
\label{sec:domain-appointments}

An \textbf{Appointment} is the central transactional entity of the
system. It links a patient, a doctor, and a specific time slot, and it
is the anchor record to which prescriptions are attached. Every
appointment passes through a defined set of lifecycle states.

\begin{figure}[h!]
  \centering
  \begin{tcolorbox}[
    enhanced, arc=4pt, boxrule=1pt,
    colback=LightGray, colframe=RuleGray,
    left=12pt, right=12pt, top=10pt, bottom=10pt,
    width=0.92\linewidth,
  ]
  \centering
  \textcolor{PrimaryBlue}{\textbf{PENDING}}
  \quad\textcolor{AccentTeal}{$\longrightarrow$}\quad
  \textcolor{MethodPost}{\textbf{CONFIRMED}}
  \quad\textcolor{AccentTeal}{$\longrightarrow$}\quad
  \textcolor{MethodGet}{\textbf{COMPLETED}}

  \vspace{8pt}

  \textcolor{PrimaryBlue}{\textbf{PENDING}} \;or\;
  \textcolor{MethodPost}{\textbf{CONFIRMED}}
  \quad\textcolor{red!70!black}{$\longrightarrow$}\quad
  \textcolor{MethodDelete}{\textbf{CANCELLED}}
  \end{tcolorbox}
  \caption{Valid appointment status transitions.}
  \label{fig:appt-lifecycle}
\end{figure}

\textbf{PENDING} is the initial state assigned at booking time. It
indicates that the appointment has been recorded in the system but the
doctor has not yet acknowledged it.

\textbf{CONFIRMED} is set by the Doctor via the
\texttt{PATCH /api/\apiversion{}/doctor/appointments/\{id\}/status}
endpoint. It indicates that the doctor has reviewed the appointment and
will attend.

\textbf{COMPLETED} is set by the Doctor after the consultation has
taken place. Once an appointment is in the COMPLETED state, it is
immutable: it cannot be rescheduled and the system will throw an
\phpclass{InvalidAppointmentStateException} if a reschedule is
attempted.

\textbf{CANCELLED} is set by the Assistant via the
\texttt{DELETE /api/\apiversion{}/assistant/appointments/\{id\}}
endpoint. Cancellation does not delete the appointment record; it
updates the status field and releases the time slot back into the
available pool by clearing the slot-availability cache key. A cancelled
appointment cannot have a prescription written against it; the system
throws an \phpclass{InvalidAppointmentStateException} if this is
attempted.

An appointment record stores the following fields:

\begin{itemize}
  \item \textbf{Doctor ID} --- the UUID of the assigned doctor.
  \item \textbf{Patient ID} --- the UUID of the patient.
  \item \textbf{Clinic ID} --- the UUID of the clinic.
  \item \textbf{Booked By} --- the UUID of the User (Assistant) who
    created the appointment record.
  \item \textbf{Appointment Date} --- the calendar date of the
    appointment in \texttt{YYYY-MM-DD} format.
  \item \textbf{Start Time} and \textbf{End Time} --- the slot's
    boundaries. The end time is always computed by the service layer
    (\texttt{start\_time + slot\_duration}); it is never accepted from
    the client.
  \item \textbf{Status} --- one of the four states above, stored as
    the enum's string value.
  \item \textbf{Notes} --- optional free-text notes provided at booking
    time.
\end{itemize}

%% ── §6.6 ─────────────────────────────────────────────────────────────────
\subsection{Prescription Management}
\label{sec:domain-prescriptions}

A \textbf{Prescription} is authored by a Doctor after a consultation
and is permanently linked to the appointment from which it originated.
The data model separates the prescription header (diagnosis, overall
notes) from the prescription line items (individual medicines).

A \textbf{PrescriptionItem} represents a single medicine entry within a
prescription. Each item records:

\begin{itemize}
  \item \textbf{Medicine Name} --- the name of the drug or compound.
  \item \textbf{Dosage} --- the amount per administration (e.g.\
    \texttt{"1 tablet"}, \texttt{"5ml"}).
  \item \textbf{Frequency} --- how often the medicine is taken (e.g.\
    \texttt{"Three times daily"}).
  \item \textbf{Duration} --- how long the course lasts (e.g.\
    \texttt{"7 days"}).
  \item \textbf{Notes} --- optional additional instructions (e.g.\
    \texttt{"Take after meals"}).
\end{itemize}

A prescription and all of its items are created in a single API call.
The \texttt{PrescriptionService::create()} method validates that the
authenticated doctor is the same doctor recorded on the appointment,
then persists the prescription header and iterates over the items array
to persist each \texttt{PrescriptionItem} record in a database
transaction to ensure atomicity.

The following constraints are enforced at the service layer:

\begin{enumerate}
  \item Only the doctor assigned to the appointment may author its
    prescription. Any other doctor attempting to do so will receive a
    \stforbid{} via \phpclass{ClinicScopeViolationException}.
  \item A prescription cannot be written for a
    \texttt{CANCELLED} appointment. Attempting to do so raises an
    \phpclass{InvalidAppointmentStateException}.
  \item A prescription cannot be viewed or updated by any doctor other
    than its author. The \texttt{PrescriptionController} guards both
    the \texttt{show} and \texttt{update} actions with an explicit
    ownership check.
\end{enumerate}

%% ─────────────────────────────────────────────────────────────────────────
\section{Clinic Scope Enforcement}
\label{sec:clinic-scope}
%% ─────────────────────────────────────────────────────────────────────────

Clinic scope enforcement is the mechanism by which the system guarantees
that a user operating within Clinic~A can never read or modify data
belonging to Clinic~B. This is implemented in two complementary layers.

\textbf{Layer 1: Middleware Injection.}
The \phpclass{ClinicScopeMiddleware} runs on every Doctor and Assistant
route. It reads the authenticated user's \texttt{clinic\_id} attribute
and calls \texttt{\$request->merge(['clinic\_id' => \$user->clinic\_id])}
to inject this value into the request's data bag. This overwrites any
\texttt{clinic\_id} value that the client may have included in the JSON
request body. The net effect is that a malicious client cannot forge a
different clinic's ID in a POST or PUT body to create or update records
outside their own clinic.

\textbf{Layer 2: Explicit Ownership Checks.}
For read operations where the resource is fetched by UUID (e.g.\
\texttt{GET /api/\apiversion{}/doctor/patients/\{id\}}), the controller
compares the fetched model's \texttt{clinic\_id} against the
\texttt{clinic\_id} injected by the middleware. If they do not match,
the controller throws a \phpclass{ClinicScopeViolationException}, which
renders as a \stforbid{} response. This ensures that even if a user
correctly guesses a resource's UUID, they cannot access it if it
belongs to a different clinic.

\begin{importantbox}[Why Two Layers?]
  The middleware alone is insufficient for read operations because the
  UUID of a resource from another clinic might be known to an attacker
  (e.g.\ obtained from a data breach, a log file, or enumeration).
  The explicit ownership check at the controller layer ensures that
  knowledge of a UUID does not grant access to it.
\end{importantbox}

%% ─────────────────────────────────────────────────────────────────────────
\section{System Boundaries and Out-of-Scope Concerns}
%% ─────────────────────────────────────────────────────────────────────────

To prevent scope creep in future development, the following capabilities
are explicitly \textbf{out of scope} for the current API version.

\begin{itemize}
  \item \textbf{Patient Portal.} Patients do not have accounts and cannot
    authenticate with the API. There is no patient-facing endpoint for
    self-registration, appointment booking, or prescription retrieval.
    These capabilities would require a separate authentication grant
    type and a new role definition.
  \item \textbf{Payment Processing.} The API contains no billing,
    invoicing, or payment integration. Financial workflows are treated
    as an external concern.
  \item \textbf{Real-time Notifications.} While the
    \texttt{SendAppointmentConfirmation} job is designed to send
    notifications, the specific notification channel (SMS, push, email)
    and the notification service provider are deployment-time
    configuration concerns external to the API's business logic.
  \item \textbf{File Attachments.} No endpoint currently accepts or
    returns binary file uploads (e.g.\ medical images, scanned
    documents). The prescription model stores text-based clinical notes
    only.
  \item \textbf{Multi-Doctor Appointments.} Each appointment is assigned
    to exactly one doctor. Scenarios involving multiple attending
    physicians per appointment are not modelled in the current schema.
  \item \textbf{Insurance and Referrals.} Third-party payer integration
    and referral tracking are external to the current domain model.
\end{itemize}

%% ─────────────────────────────────────────────────────────────────────────
\section{Key Terminology and Glossary}
\label{sec:glossary}
%% ─────────────────────────────────────────────────────────────────────────

The following terms are used with precise and consistent meaning
throughout this document.

\begin{longtable}{@{}L{3.8cm} X@{}}
  \toprule
  \rowcolor{PrimaryBlue}
  \tblhdr{Term} & \tblhdr{Definition} \\
  \midrule
  \endfirsthead

  \toprule
  \rowcolor{PrimaryBlue}
  \tblhdr{Term (cont.)} & \tblhdr{Definition} \\
  \midrule
  \endhead

  \bottomrule
  \endfoot

  API
    & Application Programming Interface. In this context, the RESTful
      HTTP interface exposed by the Laravel application. \\
  \rowalt
  Bearer Token
    & The OAuth2 Personal Access Token issued by Laravel Passport.
      Presented in the \texttt{Authorization: Bearer <token>}
      HTTP request header. \\
  Clinic Scope
    & The data isolation boundary that restricts each Doctor and
      Assistant to accessing only the records that belong to their
      assigned clinic. \\
  \rowalt
  Controller
    & A PHP class that receives an HTTP request, delegates to a
      service or repository, and returns an HTTP response. Controllers
      in this application contain no business logic. \\
  DTO
    & Data Transfer Object. A simple PHP object (or array) used to
      carry structured data between layers without Eloquent model
      coupling. \\
  \rowalt
  Eloquent
    & Laravel's ORM (Object-Relational Mapper). Each database table
      has a corresponding PHP class (a \emph{model}) that extends
      \texttt{Illuminate\textbackslash{}Database\textbackslash{}Eloquent\textbackslash{}Model}. \\
  Factory
    & A PHP class extending \texttt{Illuminate\textbackslash{}Database\textbackslash{}Eloquent\textbackslash{}Factories\textbackslash{}Factory}
      that generates realistic model instances for use in database
      seeding and PHPUnit tests. \\
  \rowalt
  Form Request
    & A dedicated Laravel class (\texttt{Illuminate\textbackslash{}Foundation\textbackslash{}Http\textbackslash{}FormRequest})
      that encapsulates validation rules and authorisation logic for
      a specific API endpoint. \\
  Guard
    & A Laravel authentication driver. The \texttt{api} guard is used
      throughout this application and is backed by Laravel Passport. \\
  \rowalt
  Happy Path
    & The test scenario in which all inputs are valid, all resources
      exist, and the response is a \stok{} or \stcreated{}. \\
  Middleware
    & A PHP class that intercepts an HTTP request before it reaches
      the controller (or an HTTP response before it leaves the
      application) to perform cross-cutting concerns such as
      authentication and scope enforcement. \\
  \rowalt
  N+1 Problem
    & A database performance anti-pattern in which loading a
      collection of~$N$ records triggers $N$~additional queries to
      load each record's relationships individually, instead of a
      single JOIN or eager-load query. \\
  Passport
    & Laravel's official OAuth2 server implementation. Provides the
      \texttt{auth:api} guard and the \texttt{createToken()} method
      on the user model. \\
  \rowalt
  Personal Access Token
    & An OAuth2 access token issued without an intermediate
      authorisation code step. Appropriate for first-party clients
      such as the clinic management dashboard. \\
  RefreshDatabase
    & A PHPUnit trait provided by Laravel that wraps each test in a
      database transaction, rolling it back after the test completes
      to ensure test isolation. \\
  \rowalt
  Repository
    & A PHP class that encapsulates all database query logic for a
      specific Eloquent model. Controllers and services depend on
      repositories, not directly on Eloquent models. \\
  Resource
    & A Laravel API resource class
      (\texttt{Illuminate\textbackslash{}Http\textbackslash{}Resources\textbackslash{}Json\textbackslash{}JsonResource})
      that transforms an Eloquent model into a JSON-serialisable
      array with explicit field control. \\
  \rowalt
  Route Model Binding
    & A Laravel feature that automatically resolves a UUID path
      segment (e.g.\ \texttt{\{id\}}) into an Eloquent model instance
      by calling \texttt{findOrFail()}, throwing a 404 if not found. \\
  Service
    & A PHP class that encapsulates business logic for a specific
      domain concept. Services are called by controllers and call
      repositories. They are the correct location for business rules
      and domain exception throwing. \\
  \rowalt
  Slot
    & A specific start-time string (format \texttt{HH:MM}) that
      represents a bookable appointment window within a doctor's
      schedule for a given day. \\
  UUID
    & Universally Unique Identifier (RFC 4122). A 128-bit value
      formatted as 32 hexadecimal digits separated by hyphens
      (\texttt{xxxxxxxx-xxxx-xxxx-xxxx-xxxxxxxxxxxx}). Used as the
      primary key for all entities in this system. \\
\end{longtable}

%% ─────────────────────────────────────────────────────────────────────────
\section{How to Use the API: A Narrative Walkthrough}
\label{sec:walkthrough}
%% ─────────────────────────────────────────────────────────────────────────

The following narrative describes the complete, end-to-end lifecycle of
a patient's clinic interaction as mediated by the API. This walkthrough
is intended to provide integrators with a concrete mental model of the
system before they consult the detailed endpoint catalogue in
Chapter~\ref{chap:endpoints}.

\textbf{Step 1 --- Clinic Setup (Super Administrator).}
A Super Administrator authenticates via
\texttt{POST /api/\apiversion{}/auth/login} and receives a bearer token.
They create a new clinic record via
\texttt{POST /api/\apiversion{}/admin/clinics}, specifying the name,
email, address, and phone number. They then create user accounts for the
clinic's medical staff via \texttt{POST /api/\apiversion{}/admin/users},
assigning each user a role (\texttt{doctor} or \texttt{assistant}) and
binding them to the clinic with the \texttt{clinic\_id} field.

\textbf{Step 2 --- Schedule Configuration (Doctor).}
The newly created Doctor authenticates and calls
\texttt{POST /api/\apiversion{}/doctor/schedules} once for each working
day, specifying the day of the week, their start time, end time, and
preferred slot duration.

\textbf{Step 3 --- Patient Registration (Assistant).}
The Assistant authenticates and registers a new patient via
\texttt{POST /api/\apiversion{}/assistant/patients}. The
\texttt{clinic\_id} is injected automatically by the middleware; the
Assistant does not need to supply it.

\textbf{Step 4 --- Slot Enquiry (Assistant).}
Before booking, the Assistant queries available slots via
\texttt{GET /api/\apiversion{}/assistant/available-slots?doctor\_id=\{id\}\&date=\{YYYY-MM-DD\}}.
The API returns the list of available start times based on the doctor's
schedule for that day, minus any already-booked slots.

\textbf{Step 5 --- Appointment Booking (Assistant).}
The Assistant books an appointment via
\texttt{POST /api/\apiversion{}/assistant/appointments}, providing the
doctor's ID, the patient's ID, the desired date, and the chosen start
time. The end time is computed automatically. A \stcreated{} response
confirms the booking. The system dispatches an asynchronous job to send
a confirmation notification.

\textbf{Step 6 --- Appointment Confirmation (Doctor).}
The Doctor sees the pending appointment in their appointment list and
calls
\texttt{PATCH /api/\apiversion{}/doctor/appointments/\{id\}/status}
with \texttt{\{"status": "confirmed"\}} to confirm their attendance.

\textbf{Step 7 --- Consultation and Prescription (Doctor).}
After the consultation, the Doctor marks the appointment as completed
via the same status endpoint with \texttt{\{"status": "completed"\}}.
They then author a prescription via
\texttt{POST /api/\apiversion{}/doctor/prescriptions}, listing the
diagnosis, any overall notes, and one or more medication items.

\textbf{Step 8 --- Prescription Retrieval.}
The prescription is permanently linked to the appointment and retrievable
at any time by the Doctor via
\texttt{GET /api/\apiversion{}/doctor/prescriptions/\{id\}}.

This eight-step workflow covers the primary use case of the system. Edge
cases, error scenarios, and the precise request/response payloads for
each step are detailed in the chapters that follow.
